% GENERATED FILE. Edit canonical lessons, then run npm run generate:books.
% canonical-source-hash: fnv1a64:f74702e15095a44f
% canonical-lessons: PA-C219-did-return, PA-C219-pita, PA-C219-pita-review, PA-C219-khadha, PA-C219-kaddhna, PA-C219-ddha, PA-C219-what-i-had, PA-C219-payoff

\chapter{I Drank the Tea}
\label{ch:pa-a2-pita}
\hlchaptermodality{219}

\begin{chapteropening}
\textbf{By the end of this chapter:} \emph{I can say what I drank, ate and took out, with a past form that agrees with the thing.}

Say what you drank and ate with pita / piti and khadha / khadhi, use kaddhna 'to take out', and recognise ddha.
\end{chapteropening}

\section[\texorpdfstring{\pa{ਮੈਂ} \pa{ਕੰਮ} \pa{ਕੀਤਾ} · \pa{ਪੀਣਾ} … (maiṁ kamm kītā …)}{ਮੈਂ ਕੰਮ ਕੀਤਾ · ਪੀਣਾ … (maiṁ kamm kītā …)}]{Return to a familiar did}

\label{lesson:PA-C219-did-return}



\begin{quote}
\emph{Hear:} \emph{maiṁ kamm kītā} Say \textbf{\pa{ਮੈਂ} \pa{ਕੰਮ} \pa{ਕੀਤਾ}}, \textquotedblleft{}I did the work\textquotedblright{}.
\end{quote}

\begin{grammarlens}[title={Grammar Lens}]
In \textbf{\pa{ਮੈਂ} \pa{ਕੰਮ} \pa{ਕੀਤਾ}}, \textbf{\pa{ਕੀਤਾ}} agrees with the work, \textbf{\pa{ਕੰਮ}}, not with the speaker: a woman says it the same way. Hold that next to two verbs that are also done to something: \textbf{\pa{ਪੀਣਾ}}, \textbf{\pa{ਖਾਣਾ}}.
\end{grammarlens}

\subsection*{Your turn}
\begin{itemize}
\raggedright
  \item \emph{Hear:} \emph{I did the work} Say \textbf{\pa{ਮੈਂ} \pa{ਕੰਮ} \pa{ਕੀਤਾ}}.
  \item \emph{Hear:} \emph{to drink} Say \textbf{\pa{ਪੀਣਾ}}.
  \item \emph{Hear:} \emph{to eat} Say \textbf{\pa{ਖਾਣਾ}}.
\end{itemize}

\begin{tcolorbox}[breakable,colback=teal!4,colframe=teal!35!black,title={Before you move on}]
Who decides the ending of \textbf{\pa{ਕੀਤਾ}}? (The work, \textbf{\pa{ਕੰਮ}}.)
\end{tcolorbox}
\section[\texorpdfstring{\pa{ਪੀਤਾ} · \pa{ਪੀਤੀ} (pītā · pītī)}{ਪੀਤਾ · ਪੀਤੀ (pītā · pītī)}]{\pa{ਪੀਤਾ}, \pa{ਪੀਤੀ} — drank}

\label{lesson:PA-C219-pita}



\begin{quote}
Say \emph{I did the work}. Then \emph{tea} and \emph{milk}.
\end{quote}

\begin{grammarlens}[title={Grammar Lens: the thing decides}]
\textbf{\pa{ਪੀਣਾ}} works like \textbf{\pa{ਕੀਤਾ}}: after \textbf{\pa{ਮੈਂ}}, \textbf{\pa{ਅਸੀਂ}} or \textbf{\pa{ਤੁਸੀਂ}}, its past agrees with what was drunk. \textbf{\pa{ਚਾਹ}} is a feminine word, \textbf{\pa{ਦੁੱਧ}} a masculine one:

\begin{quote}
\textbf{\pa{ਮੈਂ} \pa{ਚਾਹ} \pa{ਪੀਤੀ।}} — \emph{maiṁ chāh pītī.} — I drank tea.
\end{quote}

\begin{quote}
\textbf{\pa{ਮੈਂ} \pa{ਦੁੱਧ} \pa{ਪੀਤਾ।}} — \emph{maiṁ duddh pītā.} — I drank milk.
\end{quote}

A man and a woman say both lines the same way.
\end{grammarlens}

\subsection*{Your turn}
\begin{itemize}
\raggedright
  \item \emph{Hear:} \emph{we drank water} Say \textbf{\pa{ਅਸੀਂ} \pa{ਪਾਣੀ} \pa{ਪੀਤਾ।}}
  \item \emph{Hear:} \emph{you drank tea} Say \textbf{\pa{ਤੁਸੀਂ} \pa{ਚਾਹ} \pa{ਪੀਤੀ।}}
  \item \emph{Read it:} \textbf{\pa{ਮੈਂ} \pa{ਚਾਹ} \pa{ਪੀਤੀ।}} Point to \textbf{\pa{ਚਾਹ}}, then to the ending that follows it.
  \item \emph{Write it:} copy \textbf{\pa{ਪੀਤਾ}} and \textbf{\pa{ਪੀਤੀ}} once
\end{itemize}

\begin{tcolorbox}[breakable,colback=teal!4,colframe=teal!35!black,title={Before you move on}]
\emph{I drank tea}? (\textbf{\pa{ਮੈਂ} \pa{ਚਾਹ} \pa{ਪੀਤੀ}}.) \emph{I drank milk}? (\textbf{\pa{ਮੈਂ} \pa{ਦੁੱਧ} \pa{ਪੀਤਾ}}.)
\end{tcolorbox}
\section[\texorpdfstring{\pa{ਪੀਤਾ} · \pa{ਪੀਤੀ} (pītā · pītī)}{ਪੀਤਾ · ਪੀਤੀ (pītā · pītī)}]{Hear and say drank again}

\label{lesson:PA-C219-pita-review}



\begin{quote}
\emph{Hear:} \emph{pītā; pītī} Say \textbf{\pa{ਪੀਤਾ} · \pa{ਪੀਤੀ}}.
\end{quote}

\subsection*{Your turn}
\begin{itemize}
\raggedright
  \item \emph{Hear:} \emph{I drank milk} Say \textbf{\pa{ਮੈਂ} \pa{ਦੁੱਧ} \pa{ਪੀਤਾ।}}
  \item \emph{Hear:} \emph{we drank tea} Say \textbf{\pa{ਅਸੀਂ} \pa{ਚਾਹ} \pa{ਪੀਤੀ।}}
  \item \emph{Read it:} \textbf{\pa{ਤੁਸੀਂ} \pa{ਪਾਣੀ} \pa{ਪੀਤਾ।}} Give its meaning.
  \item \emph{Write it:} cover the model and write \textbf{\pa{ਮੈਂ} \pa{ਚਾਹ} \pa{ਪੀਤੀ।}} once
\end{itemize}

\begin{tcolorbox}[breakable,colback=teal!4,colframe=teal!35!black,title={Before you move on}]
Say what you drank this morning.
\end{tcolorbox}
\section[\texorpdfstring{\pa{ਖਾਧਾ} · \pa{ਖਾਧੀ} (khādhā · khādhī)}{ਖਾਧਾ · ਖਾਧੀ (khādhā · khādhī)}]{\pa{ਖਾਧਾ}, \pa{ਖਾਧੀ} — ate}

\label{lesson:PA-C219-khadha}



\begin{quote}
Say \emph{I drank tea}. Then \emph{to eat} and \emph{roti}.
\end{quote}

\begin{grammarlens}[title={Grammar Lens: a past of its own}]
\textbf{\pa{ਖਾਣਾ}} has a past of its own, \textbf{\pa{ਖਾਧਾ}}, and it agrees with the food in the same way:

\begin{quote}
\textbf{\pa{ਮੈਂ} \pa{ਰੋਟੀ} \pa{ਖਾਧੀ।}} — \emph{maiṁ roṭī khādhī.} — I ate a roti.
\end{quote}

\begin{quote}
\textbf{\pa{ਅਸੀਂ} \pa{ਸੇਬ} \pa{ਖਾਧਾ।}} — \emph{asīṁ seb khādhā.} — We ate an apple.
\end{quote}

CDIAL lists Punjabi \emph{khādhā}, \textquotedblleft{}eaten\textquotedblright{} (CDIAL 3865).
\end{grammarlens}

\subsection*{Your turn}
\begin{itemize}
\raggedright
  \item \emph{Hear:} \emph{you ate a roti} Say \textbf{\pa{ਤੁਸੀਂ} \pa{ਰੋਟੀ} \pa{ਖਾਧੀ।}}
  \item \emph{Hear:} \emph{I ate an apple and drank milk} Say \textbf{\pa{ਮੈਂ} \pa{ਸੇਬ} \pa{ਖਾਧਾ} \pa{ਅਤੇ} \pa{ਦੁੱਧ} \pa{ਪੀਤਾ।}}
  \item \emph{Read it:} \textbf{\pa{ਮੈਂ} \pa{ਰੋਟੀ} \pa{ਖਾਧੀ।}} Give its meaning.
  \item \emph{Write it:} copy \textbf{\pa{ਖਾਧਾ}} and \textbf{\pa{ਖਾਧੀ}} once
\end{itemize}

\begin{tcolorbox}[breakable,colback=teal!4,colframe=teal!35!black,title={Before you move on}]
\emph{I ate a roti}? (\textbf{\pa{ਮੈਂ} \pa{ਰੋਟੀ} \pa{ਖਾਧੀ}}.) \emph{I ate an apple}? (\textbf{\pa{ਮੈਂ} \pa{ਸੇਬ} \pa{ਖਾਧਾ}}.)
\end{tcolorbox}
\section[\texorpdfstring{\pa{ਕੱਢਣਾ} (kaḍḍhṇā)}{ਕੱਢਣਾ (kaḍḍhṇā)}]{\pa{ਕੱਢਣਾ} (kaḍḍhṇā) — to take out}

\label{lesson:PA-C219-kaddhna}



\begin{quote}
Before the new one: say \emph{I ate a roti}, then \emph{a book}.
\end{quote}

\subsection*{You'll want to know: \pa{ਕੱਢਣਾ}}
\textbf{\pa{ਕੱਢਣਾ}} — \emph{kaḍḍhṇā} — \textquotedblleft{}to take out\textquotedblright{}. It is done to something, so its past agrees with the thing:

\begin{quote}
\textbf{\pa{ਮੈਂ} \pa{ਕਿਤਾਬ} \pa{ਕੱਢੀ।}} — \emph{maiṁ kitāb kaḍḍhī.} — I took out the book.
\end{quote}

\begin{cousinweb}
\textbf{\pa{ਕੱਢਣਾ}} (\emph{kaḍḍhṇā}) is inherited: CDIAL lists Punjabi \emph{kaḍḍhṇā}, \textquotedblleft{}to pull out\textquotedblright{}, under a reconstructed \emph{*kaḍḍhati}, \textquotedblleft{}pulls, draws\textquotedblright{} (CDIAL 2660).
\end{cousinweb}

\begin{grammarlens}[title={What you've built}]
One more word for talking about the past.
\end{grammarlens}

\subsection*{Your turn}
\begin{itemize}
\raggedright
  \item \emph{Say it:} \emph{kaḍḍhṇā}
  \item \emph{Hear:} \emph{I took out the book} Say \textbf{\pa{ਮੈਂ} \pa{ਕਿਤਾਬ} \pa{ਕੱਢੀ।}}
  \item \emph{Recall:} say \emph{to eat}, \emph{to drink}, then \emph{to take out}
  \item \emph{Write it:} copy \textbf{\pa{ਕੱਢਣਾ}} once from the model
\end{itemize}

\begin{tcolorbox}[breakable,colback=teal!4,colframe=teal!35!black,title={Before you move on}]
What does \textbf{\pa{ਕੱਢਣਾ}} mean? (\textquotedblleft{}To take out\textquotedblright{}.) Say it once more.
\end{tcolorbox}
\section[\texorpdfstring{\pa{ਢ} (ḍha)}{ਢ (ḍha)}]{\pa{ਢ} — one letter from the verb to take out}

\label{lesson:PA-C219-ddha}



\begin{quote}
Say \textbf{\pa{ਕੱਢਣਾ}}, \textquotedblleft{}to take out\textquotedblright{}, the verb from the previous lesson.
\end{quote}

\subsection*{Script you'll notice: \pa{ਢ}}
\begin{quote}
\textbf{\pa{ਕ}} + \textbf{\pa{ੱ}} + \textbf{\pa{ਢ}} + \textbf{\pa{ਣ}} + \textbf{\pa{ਾ}} $\to$ \textbf{\pa{ਕੱਢਣਾ}}
\end{quote}

\begin{quote}
\textbf{\pa{ਢ}} — \emph{ḍha} — a Gurmukhi consonant
\end{quote}

Keep \textbf{\pa{ਢ}} in view. It is the third piece of \textbf{\pa{ਕੱਢਣਾ}}, after the doubling sign \textbf{\pa{ੱ}}. It belongs to the same row as \textbf{\pa{ਡ}}, which you know.

No stroke order is taught here: the track does not yet have a verified source for one.

\subsection*{Writing: copy what you can see}
Keep \textbf{\pa{ਢ}} in front of you. Trace its printed shape once with a finger, then copy it once with a pencil while the model stays visible. Say \emph{ḍha} when you finish.

\begin{quote}
This is observe-and-copy with the model visible, not writing from memory.
\end{quote}

\subsection*{Your turn}
\begin{itemize}
\raggedright
  \item \emph{Look:} at \textbf{\pa{ਢ}}; YOU SAY: \emph{ḍha}
  \item \emph{Read it:} \textbf{\pa{ਕੱਢਣਾ}}, then point to \textbf{\pa{ਢ}}
  \item \emph{Write it:} copy \textbf{\pa{ਢ}} once while the model remains visible
\end{itemize}

\begin{tcolorbox}[breakable,colback=teal!4,colframe=teal!35!black,title={Before you move on}]
Point to \textbf{\pa{ਢ}}, say \emph{ḍha}, then say the whole word \textbf{\pa{ਕੱਢਣਾ}}.

Source: \href{https://www.unicode.org/charts/PDF/U0A00.pdf}{Unicode Gurmukhi chart}.
\end{tcolorbox}
\section[\texorpdfstring{\pa{ਮੈਂ} \pa{ਰੋਟੀ} \pa{ਖਾਧੀ} \pa{ਅਤੇ} … (maiṁ roṭī khādhī …)}{ਮੈਂ ਰੋਟੀ ਖਾਧੀ ਅਤੇ … (maiṁ roṭī khādhī …)}]{\pa{ਮੈਂ} \pa{ਰੋਟੀ} \pa{ਖਾਧੀ} \pa{ਅਤੇ} \pa{ਚਾਹ} \pa{ਪੀਤੀ} — I ate a roti and drank tea}

\label{lesson:PA-C219-what-i-had}



\begin{quote}
Say \emph{I drank tea}, \emph{I ate a roti}, then \emph{I took out the book}.
\end{quote}

\begin{grammarlens}[title={Grammar Lens: two things, two endings}]
Join two pasts with \textbf{\pa{ਅਤੇ}}. Each one agrees with its own thing:

\begin{quote}
\textbf{\pa{ਮੈਂ} \pa{ਰੋਟੀ} \pa{ਖਾਧੀ} \pa{ਅਤੇ} \pa{ਚਾਹ} \pa{ਪੀਤੀ।}} — \emph{maiṁ roṭī khādhī ate chāh pītī.} — I ate a roti and drank tea.
\end{quote}

\begin{quote}
\textbf{\pa{ਅਸੀਂ} \pa{ਸੇਬ} \pa{ਖਾਧਾ} \pa{ਅਤੇ} \pa{ਦੁੱਧ} \pa{ਪੀਤਾ।}} — \emph{asīṁ seb khādhā ate duddh pītā.} — We ate an apple and drank milk.
\end{quote}
\end{grammarlens}

\subsection*{Your turn}
\begin{itemize}
\raggedright
  \item \emph{Hear:} \emph{I ate an apple and drank tea} Say \textbf{\pa{ਮੈਂ} \pa{ਸੇਬ} \pa{ਖਾਧਾ} \pa{ਅਤੇ} \pa{ਚਾਹ} \pa{ਪੀਤੀ।}}
  \item \emph{Hear:} \emph{I didn't drink milk} Say \textbf{\pa{ਮੈਂ} \pa{ਦੁੱਧ} \pa{ਨਹੀਂ} \pa{ਪੀਤਾ।}}
  \item \emph{Read it:} \textbf{\pa{ਕੱਢਣਾ}} Point to \textbf{\pa{ਢ}} and say \emph{ḍha}.
  \item \emph{Write it:} copy \textbf{\pa{ਮੈਂ} \pa{ਰੋਟੀ} \pa{ਖਾਧੀ} \pa{ਅਤੇ} \pa{ਚਾਹ} \pa{ਪੀਤੀ।}} once
\end{itemize}

\begin{tcolorbox}[breakable,colback=teal!4,colframe=teal!35!black,title={Before you move on}]
Tell your breakfast: what you ate, and what you drank.
\end{tcolorbox}
\section[\texorpdfstring{\pa{ਮੈਂ} \pa{ਚਾਹ} \pa{ਪੀਤੀ।} (maiṁ chāh pītī.)}{ਮੈਂ ਚਾਹ ਪੀਤੀ। (maiṁ chāh pītī.)}]{Breakfast: a four-skill payoff}

\label{lesson:PA-C219-payoff}



\begin{quote}
Say \emph{I drank tea}, then \emph{I ate a roti}.
\end{quote}

\subsection*{Your turn: four skills}
\begin{itemize}
\raggedright
  \item \emph{Hear:} \emph{we drank milk} Say \textbf{\pa{ਅਸੀਂ} \pa{ਦੁੱਧ} \pa{ਪੀਤਾ।}}
  \item \emph{Hear:} \emph{I took out the book} Say \textbf{\pa{ਮੈਂ} \pa{ਕਿਤਾਬ} \pa{ਕੱਢੀ।}}
  \item \emph{Read it:} \textbf{\pa{ਤੁਸੀਂ} \pa{ਰੋਟੀ} \pa{ਨਹੀਂ} \pa{ਖਾਧੀ।}} Give its meaning, then point to \textbf{\pa{ਢ}} in \textbf{\pa{ਕੱਢਣਾ}}.
  \item \emph{Write it:} cover the line and write \textbf{\pa{ਮੈਂ} \pa{ਚਾਹ} \pa{ਪੀਤੀ।}}
\end{itemize}

\begin{tcolorbox}[breakable,colback=teal!4,colframe=teal!35!black,title={Before you move on}]
Who decides \textbf{\pa{ਪੀਤਾ}} or \textbf{\pa{ਪੀਤੀ}} after \textbf{\pa{ਮੈਂ}}? (The thing drunk.)
\end{tcolorbox}
